\chapter{Conclusion}
\label{cap:conclusion}

This thesis set out to beat the existing classical baselines on speech / music classification with a causal neural model that runs in real time on a streaming input.
The project bar — macro F1 $\geq 0.85$ on a 3-class evaluation — was set up front.
Six of the seven deployed models clear it; the seventh (GMM at $0.825$) sits just below.
The neural family clears the bar by between $9$ and $13$ percentage points.

\section{Summary of contributions}
\label{sec:conclusion_contributions}

\textbf{Four neural classifiers.}
TCN (paper-faithful, $33$\,K params, macro F1 $0.9752$); TCN+LSTM (combined-experiment winner, $52$\,K, $0.9846$); SmallTCN (small-footprint, $10$\,K, $0.9766$); SmallerTCN (minimum-streaming-cost, $4.6$\,K, $0.9722$).
All four checkpoints deployed under \texttt{weights/} and reproducible from public dataset and code.

\textbf{Three classical baselines, extended to 3-class streaming.}
Lavner-and-Ruinskiy decision tree with 30-frame exponential-decay smoothing.
Khonglah-and-Prasanna GMM extended from 2-class threshold to 3-class argmax.
Khonglah-and-Prasanna SVM with a 20-frame decision-function smoothing buffer.
All three share a common feature extractor (\texttt{src/classic/feat\_extractor.py}) with parity-verified streaming equivalence.

\textbf{A 100-hour public dataset.}
\texttt{Marek324/speech-music-classification}, totalling $\sim 6\,000$ minutes across the three classes.
Built deterministically from seven public Hugging Face sources with synthetic augmentations for under-represented subclasses; Bresenham-style train/val/test split assignment guarantees every subclass appears in every split.

\textbf{Two diagnostic evaluation protocols.}
A hand-curated critical tier of $\sim 120$ adversarial clips covering known failure modes, used to surface robustness differences the standard test split obscures.
A transition analysis decomposing per-frame predictions into per-event latency and stable-region flicker at four switching cadences.

\textbf{An ablation program.}
Approximately $30$ TCN ablation variants, $8$ alternative spectral frontends, $3$ learned preprocessors, $8$ alternative backbones at two parameter budgets, $5$ combined-winner configurations, $6$ stateful tails, plus filter and stack sweeps for the SmallTCN/SmallerTCN derivation.

\textbf{A CPU complexity analysis.}
Single-threaded i5-8300H benchmark for all seven checkpoints with $5$ trials of $200$ pushes each, plus thread-scaling sweep with Amdahl-fitted parallel fractions.
Symbolic complexity metrics (params, MACs/frame, RF\_ms, $T_\infty$, $T_1 / T_\infty$) computed closed-form.

\textbf{Reproducibility infrastructure.}
Public dataset, public weights, deterministic training scripts, and three deviation documents (\texttt{differences/\{tcn,gmm\_svm,dt\}.md}) cataloguing every departure from the reference papers.
A Reflex-based web demo (\texttt{src/demo/}) runs any of the seven checkpoints on microphone or file input.

\section{Headline results}
\label{sec:conclusion_headline}

\Cref{tab:conclusion_headline} consolidates the comparison.

\begin{table}[h]
\centering
\small
\begin{tabular}{lrrrrr}
\toprule
\textbf{Model} & \textbf{Test F1} & \textbf{Crit F1} & \textbf{Median (ms)} & \textbf{RTF} & \textbf{$\Delta$rss (MB)} \\
\midrule
TCN+LSTM   & $\mathbf{0.9846}$ & $\mathbf{0.4861}$ & 11.07 & $2.1\times$ & 23 \\
TCN        & $0.9752$          & $0.4282$          & 7.18  & $3.2\times$ & 8 \\
SmallTCN   & $0.9766$          & $0.4294$          & 6.88  & $3.4\times$ & 64 \\
SmallerTCN & $0.9722$          & $0.4298$          & 4.38  & $\mathbf{5.3\times}$ & 13 \\
SVM        & $0.8750$          & $0.3763$          & 7.66  & $2.0\times$ & 9 \\
DT         & $0.8376$          & $0.3807$          & 4.15  & $2.4\times$ & 368 \\
GMM        & $0.8250$          & $0.4375$          & 3.94  & $3.8\times$ & $\mathbf{3}$ \\
\bottomrule
\end{tabular}
\caption{Consolidated headline results. Test F1 is macro F1 on the standard full-tier test split; Crit F1 is macro F1 on the adversarial critical tier; median, RTF, and $\Delta$rss are from the single-threaded i5-8300H benchmark.}
\label{tab:conclusion_headline}
\end{table}

\section{Deployment recommendations}
\label{sec:conclusion_recommendations}

The right model depends on the deployment scenario:

\begin{itemize}
    \item \textbf{Maximum F1, willing to spend the compute} $\to$ TCN+LSTM. Best on test ($0.985$), best on crit ($0.486$). Cost: $23\,\text{MB}$ marginal RSS, $2.1\times$ RTF, p99 latency $31\,\text{ms}$ (above the chunk budget; needs a jitter buffer).
    \item \textbf{Maximum F1 at minimum compute} $\to$ TCN. Test $0.975$, $8\,\text{MB}$ RSS, $3.2\times$ RTF. Reference architecture and the most defensible single choice if no specific constraint binds.
    \item \textbf{Minimum streaming compute} $\to$ SmallerTCN. Test $0.972$, $13\,\text{MB}$ RSS, $5.3\times$ RTF (highest in the family). Receptive field is $2.8\,\text{s}$ versus $8.4\,\text{s}$ for the rest of the TCN family.
    \item \textbf{Minimum memory} $\to$ GMM. Test $0.825$ (just below the project bar), crit $0.438$ (better than three of four neural variants). $3\,\text{MB}$ RSS, $3.8\times$ RTF. Footprint winner; the test-F1 cost is the price of admission. The crit-tier performance is the surprise: the GMM is more robust to adversarial conditions than three of four neural models because its asymmetric per-class scoring does not over-confidently mispredict on out-of-distribution inputs.
    \item \textbf{Adversarial robustness} $\to$ TCN+LSTM. Only checkpoint with crit-tier inactive F1 above $0.45$.
    \item \textbf{Lowest latency, deterministic CV} $\to$ DT or GMM. Median $\le 4.2\,\text{ms}$, CV $\le 10\,\%$. The DT carries a $368\,\text{MB}$ memory cost.
\end{itemize}

A consolidated reading: TCN+LSTM is the default unless a specific constraint binds; GMM is the surprise pick for memory-constrained deployments; SmallerTCN is the right pick for streaming-compute-constrained deployments; the DT for deterministic-latency-constrained deployments.
The four neural variants are interchangeable on the standard test split but differentiated on the crit tier and on the cost axes.

\section{Limitations}
\label{sec:conclusion_limitations}

\textbf{Dataset differences from the paper.}
Trained on \texttt{Marek324/speech-music-classification} ($\sim 100$\,h) rather than the paper's MUSAN+GTZAN+SSMSC+OFAI+MuSpeak+ESC+Sveriges Radio combination ($\sim 156$\,h).
Single-stage training rather than the paper's LQ pre-training plus HQ fine-tuning.
Absolute test-F1 numbers are not directly comparable to the paper's; relative ordering of variants should hold.

\textbf{Augmentation pipeline reduced to gain only.}
Only the loudness step from Schl\"uter and Grill~\cite{schluter2015augmentation} is implemented; time stretch, pitch shift, Gaussian filtering, and block mixing are skipped.
The other augmentations require operating on the pre-mel spectrum which our cache structure does not expose.

\textbf{Post-processing duration thresholds not implemented.}
The paper's §3.6 minimum-duration smoothing is not applied; we leave event-level evaluation as future work.

\textbf{Single hardware target.}
All wall-clock numbers are from one i5-8300H mobile CPU.
Server CPUs with AVX-512 would shift absolute timings; mobile CPUs without it would shift the other way.
The symbolic complexity table provides the cross-platform reference.

\textbf{Crit tier is small and hand-curated.}
$\sim 120$ clips, $\sim 30$ minutes; bootstrap CIs around $\pm 0.06$.
The clip selection is hand-curated, so the macro F1 is qualitative — ``model X tolerates this failure mode better than Y on this sample,'' not ``model X has $0.9 \pm 0.01$ on the population of all whispered speech.''

\textbf{Single seed per ablation variant.}
The combined-experiment design (\cref{sec:exp_combined}) is a partial mitigation: by stacking the source-experiment winners and confirming additivity to within $\pm 0.001$, we get evidence that the source-experiment $\Delta$s are seed-stable.
A multi-seed reproduction would close the gap.

\textbf{No 4-fold cross-validation.}
The paper uses 4-fold CV; we use a single $80/10/10$ split.
Within-thesis comparisons are valid; cross-thesis comparisons to the paper carry an additional source of variance we do not quantify.

\section{Future work}
\label{sec:conclusion_future}

\textbf{Implement the missing augmentations.}
Time stretch, pitch shift, Gaussian filtering, and block mixing.
Time stretch in particular should help the inactive-class crit-tier F1 (whisper, low-SNR speech include time-domain distortions the model has never seen).
Block mixing should help multi-speaker subclasses.

\textbf{Implement duration-threshold post-processing.}
A streaming finite-state machine that suppresses sub-threshold class transitions.
Slot into the transition-analysis harness; should improve event-level metrics without affecting frame-level macro F1.

\textbf{Multi-seed ablation reproduction.}
Three or five seeds per variant would let us make stronger claims about which $\Delta$s are seed-stable and which are seed-dependent.
Compute cost is the binding constraint: $\sim 200$ GPU-hours for a 5-seed reproduction.

\textbf{Cross-corpus evaluation.}
Evaluate on at least one external corpus (BBC Radio, Czech radio archive, GTZAN) to surface cross-language and cross-style distribution-shift effects.

\textbf{Adam+BatchNorm production checkpoint.}
The ablation program found Adam+BN tops the table at $+0.0040$ (\cref{subsec:exp_ablation}).
We did not promote it to a deployed checkpoint to keep the production recipe paper-faithful; a separate \texttt{tcn\_adam\_bn} checkpoint would let us report a higher absolute number alongside the paper-faithful one.

\textbf{Stateful-cache streaming TCN.}
The current TCN streaming implementation reprocesses the full receptive field on every push, accounting for most of the $6\times$ MAC-count gap between SmallerTCN and SmallTCN.
An idealised stateful streaming TCN would cache intermediate activations and update only the affected cells per push, dropping cost by approximately the receptive-field length.
Substantial implementation work but would close most of the streaming-cost gap.

\textbf{Retrain SVM and GMM under the latest streaming feature pipeline.}
The deployed SVM checkpoint pre-dates the latest revisions to \texttt{src/classic/feat\_extractor.py}.
A re-train on the post-revision pipeline might close some of the test-F1 gap.
